\documentclass[a4paper,12pt]{article}
\usepackage[utf8]{inputenc}
\usepackage[T1]{fontenc}
\usepackage[ngerman]{babel}
\usepackage{geometry}
\usepackage{amsmath}
\usepackage{amssymb}
\usepackage{booktabs}
\geometry{margin=2.5cm}

\begin{document}

% ================================================================
\section*{Adam: A Method for Stochastic Optimization}

\textbf{Autoren:} Diederik P. Kingma, Jimmy Lei Ba \\
\textbf{Institution:} University of Amsterdam / OpenAI; University of Toronto \\
\textbf{Konferenz:} ICLR 2015 \\
\textbf{arXiv:} 1412.6980

% ================================================================
\subsection*{1\quad Problemstellung und Motivation}

Stochastische gradientenbasierte Optimierung ist eine Grundvoraussetzung für das Training tiefer
neuronaler Netze. Klassische Verfahren wie SGD verwenden eine einheitliche Lernrate für alle
Parameter, was in hochdimensionalen, nicht-stationären Optimierungslandschaften ineffizient ist.
Zwei spezialisierte Methoden existieren als Vorgänger:
\begin{itemize}
    \item \textbf{AdaGrad} akkumuliert vergangene Gradienten und passt die Lernrate pro Parameter
    an -- wirksam bei spärlichen Gradienten, aber führt zu monoton fallender Lernrate.
    \item \textbf{RMSProp} verwendet einen exponentiellen gleitenden Durchschnitt des quadrierten
    Gradienten -- geeignet für nicht-stationäre Ziele, besitzt jedoch keine Korrektur des
    Initialisierungs-Bias.
\end{itemize}

Adam (\emph{Adaptive Moment Estimation}) kombiniert die Vorteile beider Ansätze: es schätzt
sowohl den ersten Moment (Mittelwert) als auch den zweiten Moment (unzentrierte Varianz) der
Gradienten adaptiv und korrigiert den Initialisierungs-Bias explizit.

% ================================================================
\subsection*{2\quad Methodik: Algorithmus und Update-Regel}

\paragraph{Momentschätzung.}
Zu jedem Zeitschritt $t$ wird der Gradient $g_t = \nabla_\theta f_t(\theta_{t-1})$ berechnet.
Daraus werden zwei exponentiell gewichtete gleitende Durchschnitte geschätzt:
\begin{align}
    m_t &= \beta_1 m_{t-1} + (1 - \beta_1)\, g_t, \\
    v_t &= \beta_2 v_{t-1} + (1 - \beta_2)\, g_t^2,
\end{align}
wobei $m_t$ die Schätzung des ersten Moments (Erwartungswert) und $v_t$ die Schätzung des zweiten
rohen Moments (unzentrierte Varianz) sind. Standardwerte: $\beta_1 = 0{,}9$, $\beta_2 = 0{,}999$.

\paragraph{Bias-Korrektur.}
Da $m_0 = v_0 = 0$, sind die Momentschätzungen zu Beginn des Trainings gegen null verzerrt
(besonders bei $\beta$-Werten nahe 1). Adam korrigiert diesen Bias:
\[
    \hat{m}_t = \frac{m_t}{1 - \beta_1^t}, \qquad \hat{v}_t = \frac{v_t}{1 - \beta_2^t}.
\]

\paragraph{Parameter-Update.}
Die Parameter werden gemäß
\[
    \theta_t = \theta_{t-1} - \alpha \cdot \frac{\hat{m}_t}{\sqrt{\hat{v}_t} + \epsilon}
\]
aktualisiert, wobei $\alpha$ die globale Lernrate und $\epsilon \approx 10^{-8}$ eine
Stabilisierungskonstante ist. Die effektive Schrittgröße $|\Delta_t| \lesssim \alpha$ ist
näherungsweise durch $\alpha$ beschränkt, unabhängig von der Skalierung der Gradienten.

\paragraph{Konvergenzeigenschaften.}
Adam besitzt eine $\mathcal{O}(\sqrt{T})$-Reue-Schranke für konvexe online-Optimierung, was
einer mittleren Reue von $\mathcal{O}(1/\sqrt{T}) \to 0$ für $T \to \infty$ entspricht. In
der Praxis konvergiert Adam bei nicht-konvexen Problemen (tiefe neuronale Netze) robuster und
schneller als SGD oder AdaGrad auf einer Vielzahl von Benchmarks (logistische Regression auf
MNIST, mehrschichtige Netze mit Dropout, CNNs auf CIFAR-10).

% ================================================================
\subsection*{3\quad Experimentelle Ergebnisse}

Adam wurde auf drei Problemklassen getestet:

\begin{center}
\begin{tabular}{lll}
\toprule
\textbf{Problem} & \textbf{Datensatz} & \textbf{Ergebnis} \\
\midrule
Logistische Regression & MNIST, IMDB BoW & Vergleichbar mit SGD+Nesterov, schneller als AdaGrad \\
Mehrschichtige Netze   & MNIST + Dropout  & Adam $>$ RMSProp, AdaDelta, SGD-Nesterov \\
Konvolutionale Netze   & CIFAR-10         & Adam $\approx$ SGD+Nesterov; Adagrad deutlich langsamer \\
\bottomrule
\end{tabular}
\end{center}

Darüber hinaus zeigen Ablationsstudien, dass die Bias-Korrektur bei großen $\beta_2$-Werten
(nahe 1, nötig für spärliche Gradienten) kritisch ist: Ohne Korrektur entstehen zu Beginn des
Trainings instabile, sehr große Schrittgrößen (analog zu RMSProp ohne Bias-Term).

\paragraph{AdaMax.}
Als Erweiterung leiten Kingma und Ba den AdaMax-Algorithmus her, der die $L^2$-Normierung
durch die $L^\infty$-Norm ersetzt ($u_t = \max(\beta_2 u_{t-1}, |g_t|)$). AdaMax benötigt
keine Bias-Korrektur des zweiten Moments und besitzt die einfache Schrittweiten-Schranke
$|\Delta_t| \leq \alpha$.

% ================================================================
\subsection*{4\quad Bezug zur Masterarbeit}

\paragraph{4.1\quad Verwendung als primärer Optimierer.}
In der Masterarbeit wird Adam als einziger Optimierer für das Training sowohl von U-Net als auch
FNO eingesetzt. Die Wahl ist durch seine Robustheit gegenüber der Lernrate motiviert: Da in
Experiment~1 vier verschiedene Lernraten ($10^{-3}$, $5 \times 10^{-3}$, $10^{-4}$, $5 \times 10^{-4}$)
verglichen werden, ist es wichtig, einen Optimierer zu verwenden, dessen Verhalten nicht stark
von der Wahl der Lernrate abhängt. Die adaptive Lernratenanpassung von Adam macht die
Ergebnisse vergleichbarer über verschiedene Lernraten hinweg.

\paragraph{4.2\quad Feste Lernrate ohne Scheduling.}
Im Gegensatz zu vielen Anwendungen, die Lernraten-Scheduling (z.B.\ Cosine Annealing oder
Step Decay) einsetzen, werden in der Masterarbeit feste Lernraten ohne Scheduling verwendet.
Dies entspricht der ursprünglichen Adam-Formulierung mit konstantem $\alpha$ und vereinfacht
den Hyperparameter-Vergleich in Experiment~1.

\paragraph{4.3\quad Normalisierung als Voraussetzung.}
Die in Abschnitt~4.3 der Masterarbeit beschriebene individuelle Normalisierung des Streamfunktions-Feldes
$\psi$ (Skalierung auf $\mathcal{O}(1)$-Wertebereich durch $10^{-p_\text{mean}}$) steht in
direktem Bezug zu LeCuns Empfehlungen für eingangs-normalisierte Daten. Adam reduziert zwar
die Abhängigkeit von der Input-Skalierung durch seine adaptiven Lernraten, profitiert aber
dennoch von einer guten Konditionierung des Eingaberaums -- insbesondere da die Hessian-Geometrie
des Verlustes durch die Eingabevarianz bestimmt wird.

\paragraph{4.4\quad Abgrenzung.}
Adam wurde primär für nicht-konvexe ML-Probleme mit stochastischen Mini-Batch-Gradienten entwickelt.
Die theoretischen Konvergenzeigenschaften (Reue-Schranke) gelten strikt nur für konvexe Probleme.
Für die in der Masterarbeit verwendeten tiefen U-Net- und FNO-Architekturen ist Adam empirisch
motiviert, ohne formale Konvergenzgarantien für den nicht-konvexen Fall.

\end{document}
